\documentclass[12pt,a4paper]{article}

% Packages
\usepackage[utf8]{inputenc}
\usepackage[english]{babel}
\usepackage{graphicx}
\usepackage{geometry}
\usepackage{hyperref}
\usepackage{listings}
\usepackage{xcolor}
\usepackage{float}
\usepackage{amsmath}
\usepackage{enumitem}

% Page geometry
\geometry{
    top=2.5cm,
    bottom=2.5cm,
    left=2.5cm,
    right=2.5cm
}

% Hyperref setup
\hypersetup{
    colorlinks=true,
    linkcolor=black,
    urlcolor=blue,
    citecolor=black
}

% Code listings setup
\lstset{
    language=Python,
    basicstyle=\ttfamily\small,
    keywordstyle=\color{blue},
    commentstyle=\color{green!60!black},
    stringstyle=\color{red},
    showstringspaces=false,
    breaklines=true,
    frame=single,
    numbers=left,
    numberstyle=\tiny\color{gray},
    captionpos=b,
    tabsize=4
}

\begin{document}

% Title page
\begin{titlepage}
    \centering
    \vspace*{1cm}
    
    {\Large \textbf{Instituto Superior de Engenharia de Lisboa}}\\[0.5cm]
    
    {\large Departamento de Engenharia Eletrónica e Telecomunicações e de Computadores (DEETC)}\\[1cm]
    
    {\large Mestrado em Engenharia Informática e de Computadores}\\[2cm]
    
    {\LARGE \textbf{PRACTICE 6 - Report}}\\[0.5cm]
    {\Large Train-Test Split Methods and Evaluation Metrics}\\[1.5cm]
    
    {\Large Aprendizagem e Mineração de Dados (AMD)}\\[0.3cm]
    {\large 1º Semestre -- 2025/2026}\\[3cm]
    
    \vfill
    
    {\large \textbf{Autores:}}\\[0.2cm]
    {\large
        Florin Sofinet (A46608)\\
        Diogo Guerra (A49149)\\
        Nuno Silva (A53976)
    }
    
    \vfill
    
    {\large \today}
    
\end{titlepage}

% Table of contents
\tableofcontents
\newpage

% Main content
\section{Introduction}

This report documents the implementation and analysis of exercises 6 and 7 from Practice 6, focusing on building classifiers using various train-test split methods and evaluating their performance using multiple scoring metrics. The exercises utilize the \texttt{my\_split\_and\_eval.py} script, which provides a comprehensive framework for machine learning model evaluation using scikit-learn.

\subsection{Objectives}

The main objectives of these exercises are:
\begin{itemize}
    \item Implement and understand different train-test split methodologies
    \item Build classification models using Decision Tree Classifier
    \item Apply multiple evaluation metrics to assess model performance
    \item Understand and handle metric warnings in edge cases
    \item Analyze the impact of small datasets on evaluation metrics
\end{itemize}

\section{Exercise 6: Building a Classifier with Train-Test Split}

\subsection{Overview}

Exercise 6 focuses on the fundamental implementation of a classifier using a single train-test split method. The goal is to establish the basic pipeline for training and evaluating machine learning models.

\subsection{Implementation Steps}

\subsubsection{Step a) Comment All Train-Test Split Methods}

The first step required commenting out all train-test split methods in the \texttt{list\_func\_tt\_split} list. This ensures a clean starting point where no split method is active.

\begin{lstlisting}[caption={Commented split methods}]
# list_func_tt_split = \
#   [
#    (holdout, (0.5, seed)),
#    (stratified_holdout, (1.0/3.0, seed)),
#    (repeated_holdout, (1.0/3.0, 2, seed)),
#    ...
#   ]
\end{lstlisting}

\subsubsection{Step b) Apply DecisionTreeClassifier}

The \texttt{DecisionTreeClassifier} was enabled in the \texttt{list\_func\_classifier} list:

\begin{lstlisting}[caption={Enabled Decision Tree Classifier}]
list_func_classifier = \
  [
   (DecisionTreeClassifier, ())
  ]
\end{lstlisting}

At this stage, running the script results in no output because:
\begin{enumerate}
    \item No train-test split method is active
    \item The main loop has no iterations to perform
\end{enumerate}

\subsubsection{Step c) Apply Holdout Split Method}

The holdout split method was enabled with a 50\% train-test split:

\begin{lstlisting}[caption={Enabled holdout split}]
list_func_tt_split = \
  [
   (holdout, (0.5, seed)),
  ]
\end{lstlisting}

Additionally, the \texttt{accuracy\_score} metric was enabled:

\begin{lstlisting}[caption={Enabled accuracy metric}]
list_score_metric = \
  [
   (accuracy_score, {}),
  ]
\end{lstlisting}

\subsection{Holdout Split Method}

The holdout method is the simplest train-test split technique:

\begin{itemize}
    \item \textbf{Definition:} Randomly splits the dataset into two disjoint subsets
    \item \textbf{Parameters:} 
    \begin{itemize}
        \item \texttt{test\_size=0.5} -- 50\% of data for testing
        \item \texttt{seed=5} -- Random seed for reproducibility
    \end{itemize}
    \item \textbf{Characteristics:}
    \begin{itemize}
        \item Single split execution (n\_splits=1)
        \item Fast and simple
        \item High variance in results (depends on random split)
    \end{itemize}
\end{itemize}

\subsection{Results}

Using the Pima Indians Diabetes dataset (768 instances, 8 features):

\begin{verbatim}
train_test_split: holdout
classifier: DecisionTreeClassifier
score_method: accuracy_score
::all-evaluated-datasets::
 66.15% |
66.15% (+/- 0.00%)
\end{verbatim}

\subsubsection{Analysis}

\begin{itemize}
    \item \textbf{Accuracy:} 66.15\% -- The classifier correctly predicted diabetes status in approximately 2/3 of test cases
    \item \textbf{Standard Deviation:} 0.00\% -- Expected, as only one split was performed
    \item \textbf{Interpretation:} This represents a baseline performance for the Decision Tree on this medical diagnosis task
\end{itemize}

\section{Exercise 7: Multiple Metrics and Repeated Splits}

\subsection{Overview}

Exercise 7 extends the previous work by implementing the complete scoring pipeline with multiple evaluation metrics and repeated train-test splits. This exercise also addresses important warnings that arise when working with small datasets or imbalanced class distributions.

\subsection{Implementation Steps}

\subsubsection{Step d) Complete the score\_recipe Function}

The \texttt{score\_recipe} function required two critical implementations:

\begin{lstlisting}[caption={Completed score\_recipe function}]
def score_recipe( classifier, X, y, tt_split_indexes,
                  f_score, **keyword_args_score ):
    score_all_list = list()

    for ( train_index, test_index ) in tt_split_indexes.split( X, y ):
      X_train, y_train = X[train_index], y[train_index]
      X_test,  y_test  = X[test_index],  y[test_index]

      # fit (build) model using classifier
      classifier.fit( X_train, y_train )

      # predict using the model and X_test
      y_predict = classifier.predict( X_test )

      # score the model using y_test and y_predict
      score = f_score( y_test, y_predict, **keyword_args_score )
      
      score_all_list.append( score )

    return score_all_list
\end{lstlisting}

\textbf{Key implementations:}
\begin{enumerate}
    \item \texttt{y\_predict = classifier.predict( X\_test )} -- Generates predictions for test data
    \item \texttt{score = f\_score( y\_test, y\_predict, **keyword\_args\_score )} -- Calculates the metric score
\end{enumerate}

\subsubsection{Step e) Test with Repeated Holdout}

The configuration was changed to use repeated holdout:

\begin{lstlisting}[caption={Repeated holdout configuration}]
list_func_tt_split = \
  [
   (repeated_holdout, (1.0/3.0, 2, seed)),
  ]
\end{lstlisting}

All five evaluation metrics were enabled:

\begin{lstlisting}[caption={All evaluation metrics}]
list_score_metric = \
  [
   (accuracy_score, {}),
   (precision_score, {"average":"weighted"}),
   (recall_score, {"average":"weighted"}),
   (f1_score, {"average":"weighted"}),
   (cohen_kappa_score, {}),
  ]
\end{lstlisting}

\subsection{Repeated Holdout Method}

\begin{itemize}
    \item \textbf{Definition:} Performs multiple random train-test splits
    \item \textbf{Parameters:}
    \begin{itemize}
        \item \texttt{test\_size=1.0/3.0} -- 33.33\% for testing, 66.67\% for training
        \item \texttt{n\_splits=2} -- Performs 2 independent splits
        \item \texttt{seed=5} -- Random seed for reproducibility
    \end{itemize}
    \item \textbf{Advantages:}
    \begin{itemize}
        \item Reduces variance compared to single holdout
        \item Provides standard deviation estimates
        \item More reliable performance estimation
    \end{itemize}
\end{itemize}

\subsection{Evaluation Metrics}

\subsubsection{Accuracy}

\begin{itemize}
    \item \textbf{Definition:} Ratio of correct predictions to total predictions
    \item \textbf{Formula:} $\text{Accuracy} = \frac{TP + TN}{TP + TN + FP + FN}$
    \item \textbf{Use case:} Overall correctness measure
    \item \textbf{Limitation:} Can be misleading with imbalanced datasets
\end{itemize}

\subsubsection{Precision}

\begin{itemize}
    \item \textbf{Definition:} Ratio of true positives to all positive predictions
    \item \textbf{Formula:} $\text{Precision} = \frac{TP}{TP + FP}$
    \item \textbf{Use case:} Measures the ability not to label negative samples as positive
    \item \textbf{Question answered:} "Of all instances predicted as positive, how many are actually positive?"
\end{itemize}

\subsubsection{Recall (Sensitivity)}

\begin{itemize}
    \item \textbf{Definition:} Ratio of true positives to all actual positive samples
    \item \textbf{Formula:} $\text{Recall} = \frac{TP}{TP + FN}$
    \item \textbf{Use case:} Measures the ability to find all positive samples
    \item \textbf{Question answered:} "Of all actual positive instances, how many did we correctly identify?"
\end{itemize}

\subsubsection{F1-Score}

\begin{itemize}
    \item \textbf{Definition:} Harmonic mean of precision and recall
    \item \textbf{Formula:} $\text{F1} = 2 \times \frac{\text{Precision} \times \text{Recall}}{\text{Precision} + \text{Recall}}$
    \item \textbf{Use case:} Balanced measure when you need to consider both false positives and false negatives
    \item \textbf{Range:} 0 (worst) to 1 (best)
\end{itemize}

\subsubsection{Cohen's Kappa}

\begin{itemize}
    \item \textbf{Definition:} Measures agreement between predicted and actual classifications, accounting for chance agreement
    \item \textbf{Formula:} $\kappa = \frac{p_o - p_e}{1 - p_e}$
    \begin{itemize}
        \item $p_o$ = observed agreement
        \item $p_e$ = expected agreement by chance
    \end{itemize}
    \item \textbf{Range:} -1 to +1
    \begin{itemize}
        \item $\kappa = 1$: Perfect agreement
        \item $\kappa = 0$: Agreement by chance
        \item $\kappa < 0$: Less than chance agreement
    \end{itemize}
\end{itemize}

\subsection{Results with Simple Dataset}

Using the \texttt{simple\_dataset} (6 instances, 2 features, 2 classes):

\begin{verbatim}
train_test_split: repeated_holdout
classifier: DecisionTreeClassifier

score_method: accuracy_score
::all-evaluated-datasets::
 50.00% | 0.00% |
25.00% (+/- 25.00%)

score_method: precision_score
::all-evaluated-datasets::
 100.00% | 0.00% |
50.00% (+/- 50.00%)

score_method: recall_score
::all-evaluated-datasets::
 50.00% | 0.00% |
25.00% (+/- 25.00%)

score_method: f1_score
::all-evaluated-datasets::
 66.67% | 0.00% |
33.33% (+/- 33.33%)

score_method: cohen_kappa_score
::all-evaluated-datasets::
 0.00% | 0.00% |
0.00% (+/- 0.00%)
\end{verbatim}

\subsubsection{Analysis of Results}

\begin{itemize}
    \item \textbf{Accuracy:} 25.00\% $\pm$ 25.00\% -- High variance due to small dataset
    \item \textbf{Precision:} 50.00\% $\pm$ 50.00\% -- One split had 100\%, another had 0\%
    \item \textbf{Recall:} 25.00\% $\pm$ 25.00\% -- Similar variance to accuracy
    \item \textbf{F1-Score:} 33.33\% $\pm$ 33.33\% -- Reflects the imbalance in precision/recall
    \item \textbf{Cohen's Kappa:} 0.00\% $\pm$ 0.00\% -- No better than random chance
\end{itemize}

The extremely high standard deviations and poor performance are expected with only 6 samples split into even smaller train/test sets.

\subsection{Understanding Metric Warnings}

\subsubsection{Step f) \& g) Warning Analysis}

When running the code, two types of warnings appear:

\paragraph{Warning 1: "No Predicted Samples"}

\begin{lstlisting}[language=bash, frame=single]
UndefinedMetricWarning: Precision is ill-defined and being 
set to 0.0 in labels with no predicted samples.
\end{lstlisting}

\textbf{Explanation:}
\begin{itemize}
    \item Occurs when the classifier \textbf{never predicts} a certain class
    \item Example: Test set contains both class 0 and class 1, but the classifier only predicts class 1
    \item \textbf{Formula breakdown:} $\text{Precision}_{\text{class 0}} = \frac{TP_0}{TP_0 + FP_0}$
    \item If class 0 is never predicted: $TP_0 = 0$ and $FP_0 = 0$ $\Rightarrow$ division by zero
    \item Affects: \textbf{Precision} and \textbf{F1-score}
\end{itemize}

\paragraph{Warning 2: "No True Samples"}

\begin{lstlisting}[language=bash, frame=single]
UndefinedMetricWarning: Recall is ill-defined and being 
set to 0.0 in labels with no true samples.
\end{lstlisting}

\textbf{Explanation:}
\begin{itemize}
    \item Occurs when a certain class \textbf{doesn't exist} in the test set
    \item Example: Test set only contains class 1 (no class 0 instances)
    \item \textbf{Formula breakdown:} $\text{Recall}_{\text{class 0}} = \frac{TP_0}{TP_0 + FN_0}$
    \item If no true class 0 samples exist: $TP_0 = 0$ and $FN_0 = 0$ $\Rightarrow$ division by zero
    \item Affects: \textbf{Recall} and \textbf{F1-score}
\end{itemize}

\paragraph{Why F1-Score Gets Both Warnings}

The F1-score is the harmonic mean of precision and recall:
$$\text{F1} = 2 \times \frac{\text{Precision} \times \text{Recall}}{\text{Precision} + \text{Recall}}$$

Therefore:
\begin{itemize}
    \item If precision is undefined $\Rightarrow$ F1 is undefined
    \item If recall is undefined $\Rightarrow$ F1 is undefined
    \item F1-score inherits warnings from both metrics
\end{itemize}

\subsubsection{Step h) Handling Warnings with zero\_division Parameter}

Scikit-learn provides the \texttt{zero\_division} parameter to handle these edge cases:

\begin{lstlisting}[caption={Adding zero\_division parameter}]
list_score_metric = \
  [
   (accuracy_score, {}),
   (precision_score, {"average":"weighted", "zero_division":0}),
   (recall_score, {"average":"weighted", "zero_division":0}),
   (f1_score, {"average":"weighted", "zero_division":0}),
   (cohen_kappa_score, {}),
  ]
\end{lstlisting}

\textbf{Effect of zero\_division=0:}
\begin{itemize}
    \item Sets undefined metrics to 0.0 instead of raising warnings
    \item Explicitly handles division by zero cases
    \item More appropriate than default behavior (which also sets to 0 but warns)
    \item Warnings are still generated but can be controlled separately
\end{itemize}

\subsubsection{Step i) Suppressing Warnings}

To completely suppress the warnings (with caution):

\begin{lstlisting}[caption={Suppressing UndefinedMetricWarning}]
warnings.filterwarnings( "ignore", category=UndefinedMetricWarning )
\end{lstlisting}

\textbf{Important considerations:}
\begin{itemize}
    \item \textbf{Advantage:} Cleaner output, no warning clutter
    \item \textbf{Disadvantage:} Loss of diagnostic information
    \item \textbf{Best practice:} Only use "ignore" when you fully understand the implications
    \item \textbf{Recommendation:} Keep warnings during development, suppress only in production if necessary
\end{itemize}

\section{Comparative Analysis}

\subsection{Impact of Dataset Size}

\begin{table}[H]
\centering
\begin{tabular}{|l|c|c|}
\hline
\textbf{Aspect} & \textbf{Simple Dataset (6)} & \textbf{Diabetes Dataset (768)} \\
\hline
Train size (50\%) & 3 instances & 384 instances \\
Test size (50\%) & 3 instances & 384 instances \\
Accuracy & 25\% $\pm$ 25\% & 66.15\% $\pm$ 0\% \\
Warnings & Frequent & Rare \\
Reliability & Very low & High \\
\hline
\end{tabular}
\caption{Impact of dataset size on evaluation}
\end{table}

\subsection{Holdout vs Repeated Holdout}

\begin{table}[H]
\centering
\begin{tabular}{|l|c|c|}
\hline
\textbf{Characteristic} & \textbf{Holdout} & \textbf{Repeated Holdout} \\
\hline
Number of splits & 1 & Multiple (e.g., 2, 5, 10) \\
Standard deviation & Always 0\% & Provides variance estimate \\
Reliability & Lower & Higher \\
Computation time & Faster & Slower \\
Best for & Quick testing & Reliable evaluation \\
\hline
\end{tabular}
\caption{Comparison of split methods}
\end{table}

\section{Conclusions}

\subsection{Key Learnings}

\begin{enumerate}
    \item \textbf{Train-Test Split Importance:}
    \begin{itemize}
        \item Proper data splitting is crucial for unbiased evaluation
        \item Single holdout provides quick results but high variance
        \item Repeated splits offer more reliable performance estimates
    \end{itemize}

    \item \textbf{Multiple Metrics Matter:}
    \begin{itemize}
        \item Accuracy alone is insufficient for comprehensive evaluation
        \item Precision and recall provide complementary insights
        \item F1-score balances both metrics
        \item Cohen's Kappa accounts for chance agreement
    \end{itemize}

    \item \textbf{Edge Cases and Warnings:}
    \begin{itemize}
        \item Small datasets frequently trigger metric warnings
        \item Understanding "no predicted samples" vs "no true samples" is essential
        \item The \texttt{zero\_division} parameter provides explicit control
        \item Warning suppression should be used judiciously
    \end{itemize}

    \item \textbf{Dataset Size Impact:}
    \begin{itemize}
        \item Minimum viable dataset size is crucial for reliable evaluation
        \item With 6 samples, results are essentially random
        \item Larger datasets (768+) provide stable, meaningful metrics
    \end{itemize}
\end{enumerate}

\subsection{Best Practices}

\begin{itemize}
    \item Use repeated splits (holdout, k-fold) for reliable evaluation
    \item Report multiple metrics, not just accuracy
    \item Understand warning messages before suppressing them
    \item Use \texttt{zero\_division} parameter explicitly
    \item Ensure adequate dataset size for your evaluation strategy
    \item Set random seeds for reproducibility
    \item Consider stratified splits for imbalanced datasets
\end{itemize}

\subsection{Future Work}

Potential extensions to these exercises:
\begin{itemize}
    \item Implement k-fold cross-validation for more robust evaluation
    \item Compare performance across multiple classifiers (Naive Bayes, SVM, Random Forest)
    \item Apply stratified splits to handle class imbalance
    \item Analyze learning curves with varying dataset sizes
    \item Implement bootstrap sampling for variance estimation
\end{itemize}

\end{document}
